% ~7 min
\partframe{Part IV}{Dynamic Programming}

\begin{frame}
  \frametitle{Knapsack}
  \[ K(0) = 0 \qquad K(c) = \max\,\bigl(\{0\} \cup \{\,K(c - w_i) + v_i
     \mid 0 < w_i \le c\,\}\bigr) \]
  \begin{center}
    \begin{tikzpicture}[alt={Capacities 0 to 6 in a row. The entry for
        capacity 6 reads the entries for capacities 4 and 3, that is 6
        minus weight 2 and 6 minus weight 3, for the items of weight 2 and
        value 3 and of weight 3 and value 4. Every entry is computed once
        and read many times.},
      every node/.style = {box, minimum width = 9mm, font = \small}]
      \foreach \i in {0, ..., 6} {
        \node (k\i) at (1.3 * \i, 0) {$K(\i)$};
      }
      \node[hot, minimum width = 9mm] at (k6) {$K(6)$};
      \draw[->, accent2, thick] (k6.north) to[bend right = 40]
        node[above, draw = none, fill = none, note] {$w = 2$} (k4.north);
      \draw[->, accent2, thick] (k6.south) to[bend left = 40]
        node[below, draw = none, fill = none, note] {$w = 3$} (k3.south);
    \end{tikzpicture}
  \end{center}
  Items $(w, v)$: $(2, 3)$ and $(3, 4)$, so $K(6) = \max(K(4) + 3, K(3) + 4)
  = 9$.\par\medskip
  The naive hylomorphism recurses over the items: exponential.
  \speakernote{K(c): best value for capacity c. K(6) needs K(4) and K(3)
    (weights 2 and 3). The specification recurses over the items that fit,
    one call each, so the same K(c - w) is recomputed many times. Goal:
    recurse over capacities and look earlier results up.}
\end{frame}

\begin{frame}
  \frametitle{Dynamorphisms: Algebras with a Memory}
  \begin{center}
    \begin{tikzpicture}[alt={A table for G X = 1 + X is a list of results:
        the head is the result for c, followed by the results for c minus
        1, down to 0.},
      every node/.style = {box, minimum width = 12mm, font = \small},
      node distance = 0mm]
      \node[hot] (h) {$K(c)$};
      \node[right = of h] (a) {$K(c-1)$};
      \node[right = of a] (b) {$\cdots$};
      \node[right = of b] (c) {$K(0)$};
      \node[draw = none, fill = none, note, above = 1mm of h] {head};
      \node[draw = none, fill = none, note, above = 1mm of b] {tail};
    \end{tikzpicture}
  \end{center}
  \[
    \begin{tikzcd}[alt={The dynamorphism goes from B to the table by a
        hylomorphism, then to A by taking the head.}, column sep = 40mm]
      B \arrow[r, "\hylo{\op{cons} \circ \langle a, \op{id} \rangle}{c}"]
        \arrow[rr, bend right = 25, "\op{dyna}\,a\,c"', dashed]
        \& G_\infty\,A \arrow[r, "\op{head}"] \& A
    \end{tikzcd}
  \]
  \speakernote{Table of all results below a seed: A x G (table), a list
    newest first for G X = 1 + X. A dynamorphism is a hylomorphism into
    the table whose algebra reads it; then take the head (Hinze, Wu,
    Gibbons). In the library the table is any initial algebra of Memo G A.}
\end{frame}

\begin{frame*}
  \frametitle{Dynamorphisms in Rocq}
  \rocqblock{memo.v}
  \speakernote{Memo puts an A at every node. dyna works for any initial
    algebra T of Memo G A; dyna\_iso swaps one table for another, so the
    proofs use the free one and the code can use another. ATab: tables in
    persistent arrays, one more initial algebra. (Implicit arguments dropped on the slide.)}
\end{frame*}

\begin{frame*}
  \frametitle{Changing the Recursion}
  \rocqblock{reenum.v}
  \begin{itemize}
    \item \code{an}, \code{cn}: the naive recursion (over items).
    \item \code{e}: the recursion we want (count down capacities).
    \item \code{Hlook}: \code{an} only needs results that \code{e} reaches.
  \end{itemize}
  \speakernote{reenumerate: a hylomorphism equals a dynamorphism on another
    coalgebra, if its algebra reads its subproblems through a lookup along
    that coalgebra. look\_rec recomputes; look\_tbl reads the table. Proof:
    uniqueness of hylo, with the table of the naive function. (Implicit arguments dropped.)}
\end{frame*}

\begin{frame*}
  \frametitle{Knapsack: the Calculation}
  \rocqblock[before upper = {\fvset{fontsize = \footnotesize}}]{knap-calc.v}
  {\footnotesize Step 1's proof is left out.\par}
  \speakernote{Step 1: Derive the algebra that reads subproblems through
    a lookup down the count-down (look\_down). Step 2: reenumerate with
    ATab tables; look\_tbl\_with fast\_drop\_walk swaps the walk down the
    table for a jump into the array (O(1)); dyna\_memo shares the table.}
\end{frame*}

\begin{frame*}[vertical-alignment = top]
  \frametitle{Knapsack, Step by Step}
  \vspace*{4mm}
  \begin{onlyenv}<1>
    \fvset{lastline = 1}%
    \rocqblock{knap-steps.v}
  \end{onlyenv}
  \begin{onlyenv}<2>
    \fvset{lastline = 3}%
    \rocqblock{knap-steps.v}
  \end{onlyenv}
  \begin{onlyenv}<3>
    \fvset{lastline = 5}%
    \rocqblock{knap-steps.v}
  \end{onlyenv}
  \begin{onlyenv}<4>
    \fvset{lastline = 7}%
    \rocqblock{knap-steps.v}
  \end{onlyenv}
  \begin{onlyenv}<5->
    \fvset{lastline = 11}%
    \rocqblock{knap-steps.v}
  \end{onlyenv}
  \par
  \speakernote{The goal after each rewrite, as Rocq prints it \\
    reenumerate: from items to capacities, with a table \\
    look\_tbl\_with: lookups jump into the array, O(1) \\
    dyna\_memo: each entry built once}
\end{frame*}

\begin{frame*}[vertical-alignment = top]
  \frametitle{Knapsack: What Comes Out}
  \vspace*{5mm}
  \begin{onlyenv}<1>
    \ocamlblock[before upper = {\fvset{fontsize = \scriptsize,
      baselinestretch = .84}}]{knap-arr.ml}
  \end{onlyenv}
  \begin{onlyenv}<2->
    \ocamlblock[before upper = {\fvset{fontsize = \footnotesize}}]{knap-jump.ml}
    \par\medskip
    \code{fast\_drop}: jump straight into the right block of the array,
    $O(1)$. Knapsack is linear in the capacity.
  \end{onlyenv}
  \par
  \speakernote{Real extraction output, unedited \\
    Point at \code{push}: one entry per capacity, built bottom-up \\
    Point at \code{fast\_drop}: the lookup of K(c - w) \\
    Next step: the lookup itself, a jump into the array block, O(1)}
\end{frame*}

